\section{Related literature and the contribution boundary}\label{sec:literature}

\paragraph{Private matchmaking, screening, and heterogeneous network value.}
\citet{DamianoLi2007} study an optimising monopoly matchmaker that posts type-contingent entrance-fee schedules to sort privately informed agents into complementary random bilateral matches; \citet{DamianoLi2008} analyse matching-market competition and the potentially inefficient sorting it produces. These papers already combine \emph{fee-induced selection with complementary match value}, so neither the product of two founder abilities nor price-dependent sorting is a new mechanism here. In our model the regional private providers do not control exclusive meeting places or assortative matching: each prices a distinct local founder population whose realised cross-region partners are complementary. The relevant distinction is the evaluated allocation: our premium users and basic users share the same matching technology, while a change in admission alters the mass obtaining direct basic services. The policy comparison values that restored access and the associated regional market incidence, rather than the efficiency of sorting agents into exclusive meeting places.

\citet{Hagiu2004,Hagiu2004b} study proprietary and open platform design and developers' heterogeneous contributions, while \citet{VeigaWeylWhite2017} emphasise quality-weighted sorting and the differing network values of participants. Our current uniform upper-tail selection obeys \emph{exactly} $q_i=m_i-m_i^2/2$ in every service state, and therefore can be represented by a single nonlinear participation index. We do not claim the innovation of an irreducible second network state, independent mean quality, or the general quality--quantity trade-off. \citet{PeitzSato2026} further integrate platform fees, endogenous participation and entry, and \citet{LuMatsushimaShekhar2025} connect network externalities to competitive platform pricing. The specific posted-fee response in our model is a specialisation of familiar pricing and network feedback.

\paragraph{Public alternatives and endogenous private-market responses.}
\citet{HuaChenZhang2016} consider free public and paid private services under delay- and capacity-sensitive provision and fiscal policy. \citet{LaineMa2017} derive mixed-service allocation and endogenous quality conditions; their covered-population sorting result does not become the present private-only exclusion comparison merely by setting the lower-quality price to zero. \citet{Kim2024} also studies mixed public/private market competition. \citet{AtalCuestaGonzalezOtero2024} document the entry of public pharmacies and the resulting segmentation of private customers and increases in private drug prices. The public-private payment-platform model of \citet{LiuReshidiRivadeneyra2026} demonstrates private merchant-fee increases in the presence of a welfare-oriented public platform, through differentiated preferences and network fragmentation. Therefore public entry, private fee responses, adverse distributional effects and mixed provision are all known. Our public service has no rationing, a private premium does not directly improve matching, and the distinguishing policy object is \emph{regional venture-commercialisation incidence} interacting with a fee-induced basic-access exclusion margin. In the inspected zero-fee welfare comparison of \citet{LiuReshidiRivadeneyra2026}, transaction surplus is positive; identifying that surplus with the present negative regional coefficient would violate its maintained domain. Their public-objective weights also differ from a founder profit share that enters private willingness to pay. Similarly, the logit platform-share demand in \citet{PeitzSato2026} cannot be replaced by ability-dependent upper-tail admission without changing its equilibrium restrictions. These results establish relevant mechanisms, but their particular welfare comparisons do not directly supply the public-cost derivative studied here.

\paragraph{Admission, externalities and public options.}
The general tension between private admission and external effects also precedes this analysis. \citet{KimLiXu2025} show that an entry-fee scheme can induce underparticipation even in a market with congestion, whereas nonexclusive priority search can align the revenue-maximising participation level with efficiency. Their Proposition~2 therefore rules out treating private over-exclusion in the presence of harmful participation as a new general insight here. Our comparison concerns a different margin: a technical change alters the privately selected pool, and the government values a fixed basic-plus-premium alternative with regionally assigned commercialisation returns. The premium does not improve matching priority or technological success.

\citet{OstrizekSartori2023} study screening when agents' actions generate externalities and characterise optimal contracts. A one-dimensional participation representation is therefore not a new screening class. In their February 2023 author manuscript, the principal chooses a contract menu and evaluates its externality objective; restricting that choice to our single private posted fee and introducing a separate evaluator of a predetermined access alternative does not preserve the optimal-contract conclusion. The present contribution is the explicit policy-value boundary under these restricted decisions, not the general possibility of screening with externalities.

\paragraph{Innovation policy, local market incidence and governmental objectives.}
\citet{TakaloTanayamaToivanen2013} analyse strategic interaction among R\&D project applicants, public funding agencies and private financiers with fixed project costs and financing frictions; they demonstrate that subsidy additionality and welfare can have a nontrivial relation. Thus the observation that an innovation-support programme's value cannot be judged by the count of supported projects is not itself new. Our public instrument is \emph{free basic intermediation} rather than project finance or R\&D subsidy, and premium payment does not raise the underlying research technology. \citet{DAspremontJacquemin1988,GoyalJoshi2003} study research cooperation and rivalry, while \citet{BerliantFujita2008,BerliantFujita2011} address heterogeneous knowledge. \citet{Lee2026} links local knowledge spillovers and oligopolistic rivalry to sign-dependent R\&D intervention. These studies establish that commercialisation benefits need not be positive for every market participant or region. What is specific here is that the sign of local venture-market incidence $K_i$ may be \emph{opposite} to the sign of the derivative of the local government's willingness to restore founder access with respect to technological complementarity.

\citet{LongWong2009} model two city governments and two private service firms in inter-city competition; the local-government/private-provider architecture alone is not new. Local and national investment incentives have long differed when cross-jurisdictional public-good benefits are not internalised \citep{CremerMarchandPestieau1997,MankiwWhinston1986}. Accordingly, Lemmas~1--3 are preparatory accounting/selection results, not priority claims. Nor is proof-oriented exact certification of the chosen equilibrium calibration a new theory of competition, although it makes the model's examples more credible.

\paragraph{Economic heterogeneity and policy boundaries.}
\citet{BenassiChircoColombo2025} revisit canonical mixed-oligopoly models by linking reservation-price distributions to income and asking how distributional changes shift privatization boundaries. Their analysis shows how an established model can answer a different economic question through explicit policy thresholds, while the direction of particular comparisons depends on the specified demand and cost environment. Our heterogeneity concerns founder selection and regional commercialisation incidence, rather than household income or an endogenous privatization share. The relevant comparison is the economic content of the resulting boundary, not the mere presence of additional heterogeneous agents. Here that content is the private exclusion response and the conditional profit-share band; no income-distribution result is imported.

\paragraph{The narrower contribution.}
The economically distinctive result is Proposition~1, specifically the \emph{primitive profit-share band} \eqref{eq:screen-inversion-band}: in a region where the new venture's commercialisation incidence is negative, higher technical skill complementarity can increase the publicly affordable cost of basic access because profit-related price screening excludes firms from socially valuable direct services. The condition is derived from reoptimised private fees and the local Cournot-derived $M_i$, not the assumption of a positive local innovation externality. An independent finite-$\gamma$ interval certificate proves this sign inversion on an explicitly stated continuous primitive rectangle. The primitive band separates a lower founder-share boundary for the screening response from an upper boundary for negative local incidence. The inspected sorting, platform-entry and optimal-contract propositions do not directly establish this derivative of the specified access alternative. This is a \emph{conditional, model-specific applied-theory contribution}; it is not a new general law of public options or an unrestricted classification of regional entry.

\begin{table}[htbp]\centering\small
\caption{Closest prior mechanisms and the conditional result retained here}\label{tab:prior}
\begin{tabular}{@{}>{\raggedright\arraybackslash}p{0.30\linewidth}>{\raggedright\arraybackslash}p{0.62\linewidth}@{}}\toprule
Closest prior research & Existing mechanism and distinct policy object\\\midrule
Damiano and Li (2007, 2008); Hagiu (2004b) &
Complementary-type matching, screening through admission prices and quantity--quality trade-offs are known; we value the restoration of fee-excluded basic services using local post-innovation market surplus.\\\addlinespace[3pt]
Atal et al. (2024); Liu et al. (2026); Hua et al. (2016) &
Public provision, private fee/participation responses and mixed services are known; our policy-value derivative varies with R\&D skill complementarity and locally adverse commercialisation.\\\addlinespace[3pt]
Takalo et al. (2013); Lee (2026) &
Innovation subsidies, adverse incumbent effects and sign-sensitive externalities are known; our government chooses zero-priced access against a private posted-fee selected population, not an R\&D grant.\\\addlinespace[3pt]
Peitz and Sato (2026); Cremer et al. (1997) &
Endogenous network participation and regional government wedges are established; our contribution is a primitive-level, conditional sign conflict with an exact finite-$\gamma$ continuum witness.\\\bottomrule
\end{tabular}
\end{table}

The public-cost interval that supports the three-scenario \emph{ownership switch} remains narrow and the recorded alternative skill distribution removes it. Neither the stronger continuously certified public-value \emph{slope} nor the distinct institutional labels should be used to imply a broad full-game SPNE result, an empirically identified programme effect, or a general technology-driven change in decentralised provision.